\section{Data}
\label{sec:data}

We use the student questionnaire files for all five PISA cycles from 2009 to 2022.
The 2015/2018/2022 cycles ship as SPSS \texttt{.sav}; 2009 and 2012 are fixed-width
ASCII with SPSS control files, which we parse with a purpose-built reader. After
harmonising feature names across cycles, the Albanian longitudinal file holds
27{,}042 students; a comparison set of nine countries, Albania, North Macedonia,
Montenegro, Serbia, Bulgaria (Balkan peers), Estonia and Finland (top performers),
and Mexico and Colombia (GDP-matched), holds 323{,}121 students. Raw files are
multi-gigabyte and not redistributed; processed slices are regenerated by the
pipeline.

\paragraph{Target.}
A student is \emph{at risk} if their mathematics score falls below the exact OECD
Level~2 boundary of 420.07 points. Because each score is ten plausible values, the
model-comparison target is the point (PV-mean) indicator, while the final estimates
use per-plausible-value targets combined by Rubin's rules.

\paragraph{Features.}
The core set is thirteen background variables with high cross-cycle availability:
ESCS (the OECD socioeconomic index), home possessions, parental education (HISCED)
and occupation (HISEI), gender, grade repetition, immigration status, school
belonging, teacher support in mathematics, math anxiety, grade relative to modal, and
home/school ICT resources. Composite and interaction features (a standardised SES
composite, a material-deficit rank, digital-readiness indices, and centred
interaction terms) are derived \emph{inside} the cross-validation pipeline on the
training fold only, so their scaling and rank statistics never leak from the test
fold.

\paragraph{Design-based estimates.}
The 2015/2018/2022 cycles carry 80 Fay balanced-repeated-replication weights; we
combine the BRR sampling variance with the between-plausible-value variance to report
the standard error the OECD manual prescribes. The 2009/2012 fixed-width cycles do
not carry replicate weights and fall back to an imputation-only standard error, which
we flag. Table~\ref{tab:trajectory} gives Albania's weighted at-risk trajectory with
these design-based standard errors: a V-shape that falls from 68\% (2009) to 42\%
(2018) and then snaps back to 74\% in 2022 (Figure~\ref{fig:trajectory}).

\begin{table}[t]
\centering
\caption{Albania weighted at-risk mathematics rate (below Level~2) with
design-based standard errors. BRR available only for the SPSS cycles.}
\label{tab:trajectory}
\begin{tabular}{lcccc}
\toprule
Cycle & At-risk & SE & 95\% CI & BRR \\
\midrule
2009 & 0.677 & 0.0076 & 0.663--0.692 & n/a \\
2012 & 0.607 & 0.0051 & 0.597--0.617 & n/a \\
2015 & 0.533 & 0.0108 & 0.512--0.554 & \checkmark \\
2018 & 0.424 & 0.0088 & 0.407--0.441 & \checkmark \\
\textbf{2022} & \textbf{0.739} & \textbf{0.0041} & \textbf{0.731--0.747} & \checkmark \\
\bottomrule
\end{tabular}
\end{table}

\begin{figure}[t]
\centering
\includegraphics[width=0.8\linewidth,height=0.33\textheight,keepaspectratio]{A1_atrisk_trajectory}
\caption{Albania's weighted below-Level-2 mathematics rate, 2009--2022. The 2018 low
(42\%) and the 2022 spike (74\%) bracket the structural break this paper explains.}
\label{fig:trajectory}
\end{figure}

\paragraph{A verified data gap.}
Albania genuinely lacks the ESCS index in 2012 and 2015 (an OECD coverage gap
confirmed against the source files, not a processing bug), whereas all peer countries
have full 2015 background data. This limits Albania's \emph{own} cross-cycle
socioeconomic comparison but leaves the 2018-vs-2022 shift analysis, both cycles
fully covered, unaffected.
